\chapter{67}

\section{Samstag, 1. September}



«Haallo, haallo.»

Der Ruf klang nah und ausgelassen und holte Louis vor die Hütte.
Erwartungsvoll liess er seinen Blick nach unten gleiten. Das grelle
Mittagslicht blendete. Er hielt seine flachen Hände über die Stirn. Die
hochgestreckten Arme mussten zu Joh gehören. Der winkte und führte die
kleine Gruppe an. Louis winkte zurück, setzte sich auf die oberste
Treppenstufe vor die Eingangstür und wartete. Sie schienen es nicht eilig zu haben. Er erkannte Manuela an ihrem Haar und an ihren Bewegungen. Die lange, schmale Silhouette musste zu David gehören. Die vierte Person war offensichtlich die fremde Besucherin.

Louis blieben nur mehr wenige Minuten, bis die vier oben ankommen
würden. Er stand auf. Trotz seiner vielen Ups und Downs hatte er seine
Aufgaben erfüllt. Die Nacht war mild und traumlos vorübergegangen.
Vielleicht war die Hütte sauberer als zuvor, die Schafe waren versorgt
und das Mittagessen vorbereitet. Zu gegebener Zeit würde er die
Ofentemperatur hochdrehen, um den Gratin fertig zu backen. Mission
erfüllt. Er war zufrieden.

Bald waren ihre Gesichter in Sichtweite. Schliesslich standen sie vor
ihm. Louis stutzte. Der fragende Blick der jungen Frau irritierte ihn.
Er kniff die Augen zusammen und suchte nach einer Erklärung. Sie sah
jemandem ähnlich. Er liess sich erstmal in Johs und anschliessend in
Manuelas Arme fallen. Joh machte in Richtung David eine ausladende Armbewegung.

«Das ist David. Ihn hast du kurz unten gesehen,» und zu David gewandt,
«und das ist Ciro, unser temporärer Zusenn.»

Louis begrüsste David. Der sah jünger aus, als er ihn in Erinnerung
hatte. Und zu Louis’ Erstaunen war er eine freundliche Erscheinung. Joh trat einen Schritt nach hinten, legte seinen Arm um die Schultern der jungen Frau und trat mit ihr näher an Louis heran.

«Und das ist Ava, auch ein bisschen meine Tochter,» Joh zwinkerte ihr
zu, «sie will den Rest ihrer Auszeit bei uns verbringen.»

Ava lächelte. Nach wie vor musterte sie Louis.

«Hallo, Ciro,» sie reichte ihm die Hand und stutzte, «jetzt, jetzt hab
ich’s, du bist doch mein Retter im Zug, stimmt’s?»

Manuela und Joh warfen sich erstaunte Blicke zu.

«Ich glaub’s nicht, ja, ich dachte gleich, du erinnerst mich an
jemanden. Du trägst eine andere Frisur, Wanderkleidung, freut mich.»

Louis’ Freude war echt.

«Du bist also der vielgelobte \emph{Zusenn} und Musiker? Ich bin Ava.»

Da war er wieder, dieser sympathische Blick. Er war Louis im Zug als
erstes aufgefallen. Er trat einen Schritt zurück. Was hatten sie ihr
wohl von ihm erzählt?

«Der bin ich wohl,» Louis schaute theatralisch um sich und grinste, «ich
bin der Einzige hier oben.»

Ava nickte belustigt. Nach kurzer Stille übernahm Joh. Louis war nicht
aufgefallen, dass Manuela und David bereits ihr Gepäck abgelegt und sich
vor der Hütte auf die Bank gesetzt hatten.

«Mach dir’s auch bequem, Ava. Wir kommen gleich. Ich will mich kurz von
Ciro updaten lassen.»

Joh führte Louis ein paar Schritte in Richtung Ostseite der Hütte. Sie
waren ausser Hörweite. Louis war verwirrt. Was wollte Joh von ihm
wissen? Der grinste.

«Updaten, hm, das war ein Vorwand. Ich bin überzeugt, dass hier oben
alles in bester Ordnung ist, keine Sorge. Ich will dir vielmehr etwas zu
Ava sagen,» er schloss für einen Moment die Augen und sammelte sich,
«Ava ist das einzige Kind eines meiner besten Freunde. Sie war früher
als Mädchen immer wieder mal bei uns, ferienhalber, und hilft uns
gelegentlich während der \emph{Schafscheid} aus. Vor kurzer Zeit hat sie
ihren Freund durch einen Motorradunfall verloren. Wir haben Daniel auch
gekannt, ein toller Junge,» Joh senkte den Blick, schob seine Brille
über den Nasenrücken hoch und Louis zuckte zusammen, «es fällt ihr schwer, den
Verlust zu verarbeiten. Sie nahm sich eine Auszeit und ist nun die
nächsten Wochen bei uns. Ich habe ihren Eltern versichert, dass ich auf
sie aufpasse. Ich erwarte, dass sie zuvorkommend und freundlich
behandelt wird. David habe ich bereits informiert. Ich kann mich
bestimmt auf dich verlassen.»

Joh schaute durchdringend und aussergewöhnlich ernst. Louis zögerte.

«Ja, natürlich.»

Was dachte Joh eigentlich von ihm? Was sollte er denn schon falsch
machen im Kontakt mit Ava? Louis trug ein loses Unbehagen in die Hütte. Es blieb, bis er mit den
letzten Vorbereitungen des Mittagessens ausgelastet war. Von draussen
erreichte ihn nun auch Mazis Stimme. Sie waren komplett.
